\hypertarget{chapter-21-the-day-the-keys-fell}{%
\section{Chapter 21: The Day the Keys
Fell}\label{chapter-21-the-day-the-keys-fell}}

\begin{quote}
\itshape ``Every lock in the world was built for a thief who thinks the way we
do. The quantum thief does not.''\\
\normalfont\hfill---post-quantum migration notice
\end{quote}

\hypertarget{scene-1-the-rumor-from-the-lab}{%
\subsection{Scene 1: The Rumor From the
Lab}\label{scene-1-the-rumor-from-the-lab}}

The end of every secret in the world arrived, when it came, not as an explosion
but as a paper.

Sable brought it to the rooftop at dawn, and Elena Voss knew from her face before
she spoke. It was the face of someone who had read something that could not be
un-read. ``A machine,'' Sable said. ``Not a rumor anymore. A real one, somewhere
that doesn't publish, big enough to do the thing we always said was decades off.
It can take a public key---the ordinary kind, the kind that has guarded
\emph{everything} since before either of us was born---and work backward to the
private one. Every classical signature on earth just became a lock with its key
hanging on a string beside it.''

Elena sat with the enormity of it. Not a flaw in one system. A flaw in the
\emph{idea} every old system had shared: that some sums were easy to do and
impossibly hard to undo. The quantum machine undid them. Every wallet whose
security rested on that asymmetry---every bank, every ledger, every chain that had
not prepared---was now a house whose walls had quietly turned to paper while the
occupants slept.

``How long,'' Elena asked, ``before someone with that machine starts emptying the
old wallets?''

``Some already have,'' Sable said quietly. ``The old chains are bleeding. People
are watching coins that sat untouched for a decade simply\ldots{} walk away. The
thief that thinks differently has arrived, and the locks were all built for us.''

\hypertarget{scene-2-the-ones-who-prepared}{%
\subsection{Scene 2: The Ones Who
Prepared}\label{scene-2-the-ones-who-prepared}}

But the Sigil did not bleed, and Elena understood why with a slow, vertiginous
gratitude for a decision made years before by people who had been mocked for it.

``They never trusted the easy asymmetry,'' Wren said---older now, surfacing from
whatever quiet place she lived in, drawn out by the size of the day. ``From the
first block, the Sigil signed with the strange math. The kind built on problems the
quantum thief is no better at than we are. Everyone called it paranoid. Slower.
Bigger signatures, heavier proofs, for a threat that was always `decades away.'
They paid the cost every single day for a catastrophe that might never come.'' She
almost smiled. ``It came. And the bill they'd been paying turned out to be the
cheapest insurance in history.''

It was, Elena thought, the purest vindication of the entire creed. \emph{Verify
before you trust} had a sibling no one quoted as often: \emph{prepare before you
must.} The Sigil had refused to assume the future would be kind to its
assumptions. It had bound every binary, every signature, every identity to math
that did not fall when the old math fell. The provenance that proved \emph{who}
built each thing still held---because the keys behind it had been quantum-proof
since birth.

``So we're an ark,'' Sable said.

``We're a \emph{lifeboat},'' Elena corrected, ``and the water is full of people.''

\hypertarget{scene-3-the-migration}{%
\subsection{Scene 3: The
Migration}\label{scene-3-the-migration}}

What followed was the largest rescue Elena had ever been part of, and the most
dangerous, because a rescue is exactly when the predators come.

The old chains' refugees needed to move their value into something the quantum
thief could not open. The Sigil could take them---but every migration was a moment
of exposure, an old vulnerable key signing one last time to hand its holdings to a
new safe one, and in that instant the thief with the machine could race to forge
the same signature and steal the funds mid-flight. The lifeboat had to pull people
aboard faster than the water could take them.

And the old institutions---of course, always---saw opportunity in the panic. They
offered the refugees a ``managed migration,'' a trusted service to move the funds
safely, requiring only that people hand their old keys to the institution and trust
it to do the right thing. It was the Veil's sin offered as a rescue: surrender your
secret to a privileged few, in your most desperate hour, and pray.

``They're using the quantum break to herd everyone back into trusting them,'' Sable
said, sick with it. ``The one event that proves you can't trust anyone's promises,
and they've turned it into a sales pitch for promises.''

``Then we give people the other way,'' Elena said. ``The way that needs no trust.''

\hypertarget{scene-4-no-trust-required}{%
\subsection{Scene 4: No Trust
Required}\label{scene-4-no-trust-required}}

The Sigil built the lifeboat the only way it built anything: so that no one had to
trust the people running it.

A refugee could migrate their funds through a process where the dangerous moment---
the old key's last signature---was wrapped in a proof that bound it
\emph{atomically} to its destination. The signature could only ever move the funds
to the refugee's own new quantum-proof wallet, and nowhere else. A thief who forged
it mid-flight would find the forgery could do exactly one thing: complete the rescue
the refugee had already chosen. There was nothing to steal, because the stolen
signature was good for nothing but the one honest act. No institution. No
surrendered keys. No trust required---only a proof, checkable in ten milliseconds,
that the lifeboat could not betray the people it pulled aboard.

They came by the millions. Old wallets, dormant fortunes, the savings of people who
had never heard the word \emph{provenance} and never would, all of it ferried out
of the drowning old world into math the quantum thief could not break, by a network
with no captain that nonetheless got everyone aboard.

``It held,'' Sable said, on the rooftop, watching the migration counters climb, a
whole civilization's wealth moving to safety without a single trusted hand to grease
or break. She said it the way she always did now, and Elena loved her for the habit.
``It held \emph{this time}.''

``It held because someone paid for it years ago,'' Elena said, ``when it was
boring, and unrewarded, and everyone said the threat was decades away.'' She
watched the lifeboat fill. ``That's the part no one writes the songs about. Not the
night of the catastrophe. The thousand boring days before it, when a few stubborn
people kept paying for a wall against a flood that hadn't come yet.''

The sun came up on a world whose oldest secrets had all fallen in a single season
---and on one network, prepared and patient and unbetrayable, that had spent its
whole life getting ready for exactly this, and was now the dry ground a drowning
age was climbing onto.

A stranger on a train, newly arrived from a chain that no longer existed, opened a
wallet bound to math no machine could break, and verified the tip of their new home
in ten milliseconds before the doors closed.

They had made it. The vigil, as ever, continued---on safer ground now, but
continuing, because the thief that thinks differently is never the last thief, and
the next flood is always, already, decades away.

\hypertarget{chapter-21-notes}{%
\subsection{Chapter Notes}\label{chapter-21-notes}}

\begin{itemize}
\tightlist
\item
  \textbf{The quantum break.} A sufficiently powerful quantum computer can derive
  private keys from the public keys of classical signature schemes, breaking the
  security that most existing systems rely on. The chapter imagines the day this
  becomes real.
\item
  \textbf{Post-quantum from genesis.} The Sigil signs with post-quantum
  cryptography (isogeny/lattice-style math believed hard even for quantum
  computers) from its first block---paying the daily cost of larger signatures for
  years, which becomes the decisive advantage when the break arrives. ``Prepare
  before you must,'' the sibling of ``verify before you trust.''
\item
  \textbf{Migration is the moment of danger.} Moving value off broken chains
  requires an old vulnerable key to sign once more; an attacker with the machine
  can race that signature. The defense binds the last signature \emph{atomically}
  to the refugee's own new wallet, so a forged copy can only complete the intended
  rescue---no trust in any intermediary required.
\item
  \textbf{The institutional trap, again.} The recurring antagonist exploits the
  panic by offering ``managed migration'' in exchange for surrendered keys---the
  trust-a-privileged-few sin disguised as rescue. The trustless atomic migration is
  the answer.
\end{itemize}
